\chapter{First-Order Differential Equations}

A differential equation is a sentence written in the language of change. Up to now
we have mostly started with a function and \emph{asked} for its derivative. A
differential equation turns the question around: it hands us a relationship that
the derivative must satisfy, and asks us to find the function. It says, in effect,
``I will not tell you what $y$ is. I will only tell you how fast $y$ changes, and I
will let that rate depend on where you are and what $y$ already is. Now go find
$y$.''

This is not a niche trick. It is the way almost all of quantitative science is
actually written. Newton did not hand us the position of a planet; he handed us a
rule about its acceleration. A biologist does not start with a population curve;
she starts with a statement that the growth rate is proportional to the population.
A bank does not tell you your balance in ten years; it tells you the interest rate,
which is a statement about $\dv{B}{t}$. The natural laws of the world are
overwhelmingly written as differential equations, and the act of \emph{solving}
one is the act of turning a local rule about change into a global picture of
behavior.

A \textbf{first-order} differential equation is one in which only the first
derivative appears: no $y''$, no $y'''$, just $y'$ and $y$ and the independent
variable. The general shape is
\[
  \dv{y}{x} = f(x, y).
\]
Read it slowly. At every point $(x,y)$ in the plane, the right-hand side $f(x,y)$
hands you a \emph{number} --- a slope. So a first-order ODE is really a field of
little slope-marks blanketing the plane, and a solution is a curve that is
everywhere tangent to those marks. That picture, the \emph{slope field}, is the
single most important image in this chapter. Hold onto it.

\section{The Picture First: Slope Fields}

Before we solve anything, let us \emph{see} what an equation like
$\dv{y}{x} = f(x,y)$ is asking. The equation assigns a slope to every point. If we
draw a short line segment of that slope at a grid of points, we get a slope field
(also called a direction field). A solution curve is then nothing more than a path
that flows along the segments, like a leaf carried on a current.

\begin{intuitionbox}
Think of the slope field as a current in a river and the solution curves as the
paths leaves take when dropped in. Drop a leaf at a point --- that is choosing an
initial condition --- and the current carries it along a unique path. Different
starting points give different paths, but they all obey the same current. Solving
the ODE means \emph{finding the shape of those paths}. You can often sketch their
shape by eye, long before you can write a formula, just by following the arrows.
\end{intuitionbox}

Consider the simplest interesting example, $\dv{y}{x} = y$. At a point with height
$y=2$ the slope is $2$; at height $y=-1$ the slope is $-1$; along the $x$-axis the
slope is $0$. The slope depends only on height, not on $x$, so the field looks the
same as you slide left or right. Figure~\ref{ch37:slopefield-exp} shows it.

\begin{figure}[H]
\centering
\begin{tikzpicture}
\begin{axis}[
  width=11cm, height=7cm,
  axis lines=middle,
  xlabel=$x$, ylabel=$y$,
  xmin=-2.2, xmax=2.2, ymin=-2.6, ymax=2.6,
  xtick={-2,-1,0,1,2}, ytick={-2,-1,0,1,2},
  enlargelimits=false,
]
% slope field for y' = y : at each grid point draw a small segment of slope y.
% We normalize each segment to a fixed visual length so steep marks stay tidy.
\foreach \xx in {-2,-1.5,...,2}{
  \foreach \yy in {-2,-1.5,...,2}{
    \pgfmathsetmacro{\nrm}{0.20/sqrt(1 + (\yy)*(\yy))}%
    \edef\temp{\noexpand\draw[gray, line width=0.5pt]
      (axis cs:\xx-\nrm, \yy-\nrm*\yy) -- (axis cs:\xx+\nrm, \yy+\nrm*\yy);}
    \temp
  }
}
% a few solution curves y = C e^x
\addplot[curveblue, line width=1.3pt, domain=-2.2:1.0, samples=80]
  {0.35*exp(x)};
\addplot[curveblue, line width=1.3pt, domain=-2.2:1.6, samples=80]
  {-0.35*exp(x)};
\addplot[tangentred, line width=1.3pt, domain=-2.2:0.55, samples=80]
  {0.9*exp(x)};
\end{axis}
\end{tikzpicture}
\caption{The slope field of $\dv{y}{x}=y$. Each gray segment has slope equal to its
own height (so the marks are flat along the $x$-axis and steepen away from it).
Three solution curves $y=Ce^{x}$ (blue and red) thread through the field, everywhere
tangent to it. Notice how the curves never cross --- one leaf, one path.}
\label{ch37:slopefield-exp}
\end{figure}

The blue and red curves are the solutions $y = Ce^{x}$. We have not derived that
yet, but you can \emph{believe} it from the picture: a curve hugging arrows whose
slope equals their height is exactly a curve whose own slope equals its own height,
and the exponential is the function that is its own derivative.

\begin{happeningbox}
A first-order ODE is a \emph{local} law: it only ever talks about one point at a
time, telling you the slope there. Solving it means \emph{integrating that
local law into a global shape}. The remarkable thing --- guaranteed by the
Fundamental Theorem of Calculus and the existence--uniqueness theorems --- is that
an infinitesimal rule, repeated, knits itself into a single coherent curve. Every
solution method in this chapter is a different way of doing that knitting.
\end{happeningbox}

\begin{focusbox}
Three words you must keep straight:
\begin{itemize}
\item A \textbf{general solution} is a whole \emph{family} of curves, carrying one
arbitrary constant $C$ (because a first-order equation involves one integration,
and integration always leaves a $+C$).
\item An \textbf{initial condition} like $y(x_0)=y_0$ picks \emph{one} curve out of
the family by naming a point it must pass through. This nails down $C$.
\item A \textbf{particular solution} is that single chosen curve.
\end{itemize}
The general solution is the river; the initial condition is where you drop the
leaf; the particular solution is the leaf's path.
\end{focusbox}

\section{Separable Equations}

The friendliest equations are those in which the slope $f(x,y)$ factors into a part
that depends only on $x$ times a part that depends only on $y$:
\[
  \dv{y}{x} = g(x)\,h(y).
\]
We call these \textbf{separable}, because we can herd all the $y$'s to one side and
all the $x$'s to the other. The recipe taught in most courses is breathtakingly
casual: treat $\dv{y}{x}$ as a fraction, multiply both sides by $\dd x$, divide by
$h(y)$, and integrate:
\[
  \frac{\dd y}{h(y)} = g(x)\,\dd x
  \quad\Longrightarrow\quad
  \int \frac{\dd y}{h(y)} = \int g(x)\,\dd x.
\]
This works. But notice what we just did: we chopped the derivative $\dv{y}{x}$ in
half and treated $\dd y$ and $\dd x$ as separate movable objects, as if the
derivative were a genuine ratio of two numbers. It is not. $\dv{y}{x}$ is a
\emph{limit} of ratios, a single indivisible symbol. So why is this allowed?

\begin{whybox}
\textbf{Separation of variables is the chain rule run backwards.} Here is the
rigorous derivation, with no division of $\dd y$ by $\dd x$.

Start from $\dv{y}{x}=g(x)\,h(y)$. Assume $h(y)\neq 0$, and divide both sides
(legitimately, these are real numbers) by $h(y)$:
\[
  \frac{1}{h(y)}\,\dv{y}{x} = g(x).
\]
Now introduce $H(y)$, an antiderivative of $1/h(y)$, so that $H'(y)=1/h(y)$. By the
\emph{chain rule}, the left-hand side is exactly the derivative of $H(y(x))$ with
respect to $x$:
\[
  \ddx{H(y(x))} = H'(y)\,\dv{y}{x} = \frac{1}{h(y)}\,\dv{y}{x}.
\]
So our equation says
\[
  \ddx{H(y(x))} = g(x).
\]
Two functions of $x$ with the same derivative differ by a constant, so integrating
both sides in $x$ gives
\[
  H(y) = \int g(x)\,\dd x + C.
\]
That is \emph{precisely} the answer the casual ``multiply by $\dd x$'' recipe
produces, because $\int \dd y / h(y)$ is just $H(y)$. The fraction-juggling is a
mnemonic; the chain rule is the mechanism. The $\dd y$ and $\dd x$ are not really
being separated --- they are bookkeeping symbols that happen to track the
substitution $u=y(x)$ correctly. This is the same reason $u$-substitution in
integration works, and it is no accident: separation of variables \emph{is}
$u$-substitution wearing a different hat.
\end{whybox}

\begin{trapbox}
When we divided by $h(y)$ we assumed $h(y)\neq 0$. But any value $y=c$ with
$h(c)=0$ gives the \emph{constant} solution $y\equiv c$ (its derivative is zero, and
$g(x)\,h(c)=0$ too). These constant ``equilibrium'' solutions are real solutions
that the separation recipe can quietly lose, because you divided by zero to get
them. Always check the roots of $h$ separately. We will meet these again as
equilibria in the section on autonomous equations.
\end{trapbox}

\subsection{A First Worked Example}

Solve $\dv{y}{x} = xy$ with $y(0)=3$.

\begin{walkbox}
Is it separable? The right side is $x \cdot y$, a function of $x$ times a function
of $y$. Yes, with $g(x)=x$ and $h(y)=y$. So we expect to find $\int \dd y/y$ on the
left and $\int x\,\dd x$ on the right.
\end{walkbox}

\textbf{Step 1. Separate.} Divide by $y$ (noting $y=0$ as a possible lost solution):
\[
  \frac{1}{y}\,\dv{y}{x} = x.
\]
\textbf{Step 2. Integrate both sides in $x$.} The left integrates to $\ln\abs{y}$
(chain rule in reverse), the right to $\tfrac12 x^2$:
\[
  \ln\abs{y} = \tfrac12 x^2 + C.
\]
\textbf{Step 3. Solve for $y$.} Exponentiate:
\[
  \abs{y} = e^{C}\,e^{x^2/2}, \qquad\text{so}\qquad y = A\,e^{x^2/2},
\]
where $A=\pm e^{C}$ is a new constant absorbing the sign and the $e^C$. (Letting
$A=0$ even recovers the lost solution $y\equiv 0$, so the single formula
$y=Ae^{x^2/2}$ captures everything.)

\textbf{Step 4. Apply the initial condition.} $y(0)=3$ means $3 = A\,e^{0}=A$, so
$A=3$. The particular solution is
\[
  \boxed{\,y = 3\,e^{x^2/2}\,}.
\]
Check: $\dv{y}{x} = 3\,e^{x^2/2}\cdot x = x\cdot y$. It satisfies the equation, and
$y(0)=3$. Done.

\begin{labbox}
Solve $\dv{y}{x} = \dfrac{x}{y}$ with $y(0)=2$ before reading on. Separate as
$y\,\dd y = x\,\dd x$, integrate to $\tfrac12 y^2 = \tfrac12 x^2 + C$, so
$y^2 - x^2 = $ const. The solution curves are \emph{hyperbolas}. The initial
condition $y(0)=2$ gives $4 - 0 = 4$, so $y = \sqrt{x^2+4}$. Notice the answer is an
implicit relation $y^2-x^2=4$ that we then solved for $y$; many separable problems
leave you with an implicit equation, and that is a perfectly good answer.
\end{labbox}

\subsection{Why an Initial Condition Pins Down One Curve}

\begin{connectionbox}
The general solution $y=Ae^{x^2/2}$ is a one-parameter family --- an entire stack of
curves, one for each value of $A$. They fill the plane without crossing, just like
the leaves' paths in the river picture. The existence and uniqueness theorem (which
holds whenever $f(x,y)$ and $\pdv{f}{y}$ are continuous near the starting point)
guarantees that through each point passes \emph{exactly one} solution curve. That is
why naming one point, the initial condition, selects one and only one member of the
family. No two distinct solution curves of a well-behaved first-order ODE ever
touch.
\end{connectionbox}

\section{Exponential Growth and Decay}

The most important separable equation in all of science is
\[
  \dv{y}{t} = k\,y,
\]
which says \emph{the rate of change is proportional to the amount present}. Read it
as a sentence: ``the more there is, the faster it grows'' (if $k>0$), or ``the more
there is, the faster it shrinks'' (if $k<0$). Money in a continuously compounded
account, a colony of bacteria with unlimited food, a lump of radioactive uranium,
the cooling-off of a difference in temperature --- all obey this.

Separating, $\dfrac{\dd y}{y} = k\,\dd t$, so $\ln\abs{y} = kt + C$, and
\[
  \boxed{\,y(t) = y_0\,e^{kt}\,},
\]
where $y_0=y(0)$ is the starting amount. If $k>0$ the amount explodes; if $k<0$ it
decays toward zero. Figure~\ref{ch37:growth-decay} shows both.

\begin{figure}[H]
\centering
\begin{tikzpicture}
\begin{axis}[
  width=12cm, height=6.5cm,
  axis lines=middle,
  xlabel=$t$, ylabel=$y$,
  xmin=0, xmax=4, ymin=0, ymax=6,
  xtick={0,1,2,3,4}, ytick={1,2,3,4,5},
  legend pos=north west, legend cell align=left,
  domain=0:4, samples=100,
]
\addplot[curveblue, line width=1.3pt] {exp(0.6*x)};
\addlegendentry{growth, $k=0.6$}
\addplot[tangentred, line width=1.3pt] {5*exp(-0.6*x)};
\addlegendentry{decay, $k=-0.6$}
\addplot[gray, dashed, line width=0.8pt] {1};
\end{axis}
\end{tikzpicture}
\caption{Exponential growth ($k>0$) and decay ($k<0$). Both start at the dashed
reference and obey $y'=ky$. Growth curves upward without bound; decay falls toward
the $t$-axis but never reaches it.}
\label{ch37:growth-decay}
\end{figure}

\begin{historybox}
The radioactive decay law $N'=-\lambda N$ was pieced together by Ernest Rutherford
and Frederick Soddy around 1902 from careful measurements of how thorium emanation
faded. They discovered that a fixed \emph{fraction} of atoms decayed per unit time,
regardless of how many there were --- the signature of $N'\propto N$. From this came
the concept of \emph{half-life}, the time $t_{1/2}$ for half the sample to vanish.
Setting $\tfrac12 = e^{-\lambda t_{1/2}}$ gives the clean and famous relation
$t_{1/2} = \dfrac{\ln 2}{\lambda}$. Half a century later this same equation, run in
reverse, let Willard Libby date the dead from the carbon-14 left in their bones, an
idea that won the 1960 Nobel Prize in Chemistry.
\end{historybox}

\begin{applicationbox}
\textbf{Carbon dating, worked.} Living tissue holds a steady ratio of carbon-14;
when it dies, the $^{14}$C decays with half-life $t_{1/2}=5730$ years, so
$\lambda = \dfrac{\ln 2}{5730}\approx 1.21\times 10^{-4}\ \text{yr}^{-1}$. A bone is
found with $62\%$ of the living $^{14}$C level. How old is it? Solve
$0.62 = e^{-\lambda t}$ for $t$:
\[
  t = -\frac{\ln 0.62}{\lambda}
    = -\frac{\ln 0.62}{1.21\times 10^{-4}}
    \approx \frac{0.478}{1.21\times 10^{-4}}
    \approx 3950\ \text{years}.
\]
The bone is roughly four thousand years old. The whole calculation is one
separable ODE, solved once and reused forever.
\end{applicationbox}

\subsection{Newton's Law of Cooling}

A hot cup of coffee in a cool room cools fast at first, then slowly, leveling off at
room temperature. Newton's law of cooling captures this: the rate of cooling is
proportional to the temperature \emph{difference} between the object and its
surroundings. If $T(t)$ is the object's temperature and $T_{\!a}$ the (constant)
ambient temperature,
\[
  \dv{T}{t} = -k\,(T - T_{\!a}), \qquad k>0.
\]
The minus sign is the physics: if $T>T_{\!a}$ the object is hotter than the room, so
$\dv{T}{t}<0$ and it cools; if $T<T_{\!a}$ it warms. Either way it heads toward
$T_{\!a}$.

\begin{whybox}
This looks more complicated than $y'=ky$, but it is the \emph{same} equation in
disguise. Let $u = T - T_{\!a}$ be the temperature \emph{difference}. Since $T_{\!a}$
is constant, $\dv{u}{t} = \dv{T}{t}$, and the equation becomes
\[
  \dv{u}{t} = -k\,u,
\]
plain exponential decay! The difference from ambient decays exponentially:
$u(t)=u_0 e^{-kt}$. Translating back,
\[
  T(t) = T_{\!a} + (T_0 - T_{\!a})\,e^{-kt}.
\]
The substitution $u=T-T_{\!a}$ \emph{slid the whole problem down} so that the
equilibrium sat at zero, where we already know the answer, then slid it back up. A
huge amount of differential-equations technique is exactly this: find the change of
variable that turns your problem into one you have already solved.
\end{whybox}

\begin{applicationbox}
\textbf{The coffee, worked.} Coffee at $T_0 = 90^\circ$C sits in a $20^\circ$C room.
After $5$ minutes it has cooled to $70^\circ$C. When will it reach a drinkable
$50^\circ$C?

First find $k$. With $T_{\!a}=20$ and $T_0=90$, the difference starts at $70$ and at
$t=5$ has fallen to $70-20=50$:
\[
  50 = 70\,e^{-5k}\ \Rightarrow\ e^{-5k}=\tfrac{5}{7}\ \Rightarrow\
  k = \tfrac{1}{5}\ln\tfrac{7}{5}\approx 0.0673\ \text{min}^{-1}.
\]
Now solve for the time when $T=50$, i.e.\ the difference is $30$:
\[
  30 = 70\,e^{-kt}\ \Rightarrow\ e^{-kt}=\tfrac{3}{7}\ \Rightarrow\
  t = \frac{\ln(7/3)}{k}\approx \frac{0.847}{0.0673}\approx 12.6\ \text{min}.
\]
So about $12.6$ minutes after pouring. Notice the curve in
Figure~\ref{ch37:cooling} flattening toward the room line --- that flattening is the
$e^{-kt}\to 0$ tail.
\end{applicationbox}

\begin{figure}[H]
\centering
\begin{tikzpicture}
\begin{axis}[
  width=12cm, height=6cm,
  axis lines=middle,
  xlabel=$t$ (min), ylabel=$T$ ($^\circ$C),
  xmin=0, xmax=40, ymin=0, ymax=100,
  xtick={0,10,20,30,40}, ytick={20,40,60,80,100},
  domain=0:40, samples=100,
]
\addplot[curveblue, line width=1.3pt] {20 + 70*exp(-0.0673*x)};
\addplot[gray, dashed, line width=0.9pt] {20};
\node[gray] at (axis cs:32,26) {ambient $T_a=20$};
\end{axis}
\end{tikzpicture}
\caption{Newton cooling: the coffee's temperature decays toward the ambient line
$T_a=20^\circ$C. Steep at first (big difference), then ever gentler. The temperature
\emph{difference} is what decays exponentially.}
\label{ch37:cooling}
\end{figure}

\section{First-Order Linear Equations and the Integrating Factor}

Not every equation separates. Consider
\[
  \dv{y}{x} + 2y = x.
\]
Try to separate it: $\dv{y}{x} = x - 2y$. The right side is $x-2y$, which does
\emph{not} factor as (function of $x$)(function of $y$). Stuck. We need a new idea.

The class of equations we can crack is the \textbf{first-order linear} equations:
\[
  \dv{y}{x} + P(x)\,y = Q(x).
\]
``Linear'' means $y$ and $y'$ appear only to the first power, never multiplied
together, never inside a sine or a square. The coefficients $P$ and $Q$ may be any
functions of $x$. This single template covers an enormous range of models.

\subsection{The Idea: Make the Left Side a Single Derivative}

Here is the dream. The left side $\dv{y}{x} + P(x)\,y$ is \emph{almost} the
derivative of a product. Recall the product rule:
\[
  \ddx{\mu y} = \mu\,\dv{y}{x} + \mu' y.
\]
If only we could multiply our whole equation by some cleverly chosen function
$\mu(x)$ so that the left side became exactly $\ddx{\mu y}$, then the equation would
collapse to
\[
  \ddx{\mu y} = \mu\,Q,
\]
and we could just \emph{integrate both sides} and be done. The function $\mu$ that
performs this magic is called the \textbf{integrating factor}. Let us not memorize
it --- let us \emph{derive} what it has to be.

\begin{whybox}
\textbf{Deriving the integrating factor.} Multiply the equation by an unknown
positive function $\mu(x)$:
\[
  \mu\,\dv{y}{x} + \mu P\,y = \mu Q.
\]
We \emph{want} the left side to equal $\ddx{\mu y} = \mu\,\dv{y}{x} + \mu' y$.
Comparing the two expressions term by term, the $\mu\,y'$ pieces already match. The
remaining pieces match precisely when
\[
  \mu' = \mu P.
\]
But that is a \emph{separable} equation for $\mu$! Separate it:
$\dfrac{\dd\mu}{\mu} = P\,\dd x$, integrate to $\ln\mu = \int P\,\dd x$, and
exponentiate:
\[
  \boxed{\,\mu(x) = e^{\int P(x)\,\dd x}\,}.
\]
That is where the famous formula comes from. It is not pulled from a hat; it is
\emph{forced} on us by demanding that the product rule fit. The integrating factor
is exactly the function whose logarithmic derivative is $P$, because that is what it
takes to make $\mu y'+\mu P y$ collapse into one clean derivative.
\end{whybox}

\begin{happeningbox}
What is really happening is a \emph{change of viewpoint}. The combination
$y' + Py$ does not look like a derivative of anything tidy. But multiply by $\mu$
and it becomes the derivative of the single quantity $\mu y$. We have re-expressed
an awkward two-term expression as the rate of change of one product. Once something
is written as $\ddx{\text{stuff}}$, the Fundamental Theorem of Calculus lets us undo
it by integration. The integrating factor is a lens that brings the left side into
focus as a single derivative.
\end{happeningbox}

\subsection{The Full Recipe, Derived}

Putting it together: to solve $\dv{y}{x} + P(x)\,y = Q(x)$,
\begin{enumerate}
\item Compute the integrating factor $\mu = e^{\int P\,\dd x}$.
\item Multiply the equation by $\mu$. The left side is now $\ddx{\mu y}$ by
construction.
\item Integrate both sides: $\mu y = \int \mu Q\,\dd x + C$.
\item Divide by $\mu$ to get $y = \dfrac{1}{\mu}\left(\int \mu Q\,\dd x + C\right)$.
\end{enumerate}

\subsection{Worked Example with the Integrating Factor}

Solve $\dv{y}{x} + 2y = x$ with $y(0)=1$.

\begin{walkbox}
Here $P(x)=2$ (a constant) and $Q(x)=x$. The integrating factor will be
$e^{\int 2\,\dd x}=e^{2x}$. We expect to multiply through by $e^{2x}$ and watch the
left side fold into $\ddx{e^{2x}y}$.
\end{walkbox}

\textbf{Step 1. Integrating factor.} $\mu = e^{\int 2\,\dd x}=e^{2x}$.

\textbf{Step 2. Multiply through.}
\[
  e^{2x}\dv{y}{x} + 2e^{2x}y = x\,e^{2x}.
\]
The left side is exactly $\ddx{e^{2x}y}$ --- check it with the product rule:
$\ddx{e^{2x}y} = e^{2x}y' + 2e^{2x}y$.

\textbf{Step 3. Integrate both sides.}
\[
  e^{2x}y = \int x\,e^{2x}\,\dd x.
\]
The integral on the right needs integration by parts (let $u=x$, $\dd v=e^{2x}\dd
x$):
\[
  \int x\,e^{2x}\,\dd x = \frac{x}{2}e^{2x} - \frac{1}{2}\int e^{2x}\,\dd x
  = \frac{x}{2}e^{2x} - \frac{1}{4}e^{2x} + C.
\]
So
\[
  e^{2x}y = \frac{x}{2}e^{2x} - \frac{1}{4}e^{2x} + C.
\]
\textbf{Step 4. Solve for $y$.} Divide by $e^{2x}$:
\[
  y = \frac{x}{2} - \frac{1}{4} + C\,e^{-2x}.
\]
\textbf{Step 5. Initial condition.} $y(0)=1$:
\[
  1 = 0 - \tfrac14 + C \ \Rightarrow\ C = \tfrac54.
\]
The particular solution is
\[
  \boxed{\,y = \frac{x}{2} - \frac14 + \frac54\,e^{-2x}\,}.
\]
Notice the structure: a part $\frac{x}{2}-\frac14$ that tracks the forcing $Q=x$
(the ``steady'' response), plus a transient $\frac54 e^{-2x}$ that fades away. This
split into a lingering particular response and a dying homogeneous transient recurs
everywhere in applied mathematics.

\begin{trapbox}
Two ways people wreck this method:
\begin{enumerate}
\item Forgetting to put the equation in \emph{standard form} first. The coefficient
of $y'$ must be $1$ before you read off $P$. If you start from
$x\,y' + y = x^2$, divide by $x$ to get $y' + \tfrac1x y = x$, so $P=1/x$, not
something else.
\item Dragging the $+C$ from the integrating-factor integral $\int P\,\dd x$ into
$\mu$. You do not need it: any one antiderivative works, because an extra constant
in the exponent just multiplies $\mu$ by a constant, and that constant cancels when
you divide by $\mu$ at the end. Use the simplest $\mu$. The $+C$ that \emph{does}
matter is the one from Step 3, integrating $\mu Q$.
\end{enumerate}
\end{trapbox}

\begin{connectionbox}
The integrating-factor trick is the same spirit as ``completing the square'': you
massage an expression by a deliberate multiplication (or addition) until it
becomes a recognizable, invertible shape. There, you force a perfect square; here,
you force a perfect derivative. Recognizing the shape of a known operation, and
bending your problem until it fits, is one of the great recurring moves in all of
mathematics.
\end{connectionbox}

\section{Autonomous Equations, Equilibria, and Stability}

A first-order equation is \textbf{autonomous} when the right side does not mention
the independent variable explicitly:
\[
  \dv{y}{t} = f(y).
\]
The slope depends only on \emph{where} you are ($y$), never on \emph{when} ($t$).
Newton cooling, exponential growth, and the logistic equation we are about to meet
are all autonomous. Autonomy buys us a lot: it lets us read off the long-term
fate of solutions \emph{without solving the equation at all}, just by looking at the
sign of $f$.

\subsection{Equilibria: Where Nothing Changes}

An \textbf{equilibrium} (or steady state, or fixed point) is a value $y^*$ where
$f(y^*)=0$. There, $\dv{y}{t}=0$: the system is frozen. If you start exactly at
$y^*$ you stay there forever, giving the constant solution $y\equiv y^*$. These are
the very ``lost'' constant solutions from the separable section, now promoted to
center stage.

The deep question is \emph{stability}: if you nudge the system slightly off an
equilibrium, does it slide back (stable) or run away (unstable)?

\begin{intuitionbox}
Picture a marble in a landscape. A \textbf{stable} equilibrium is the bottom of a
valley: push the marble a little and it rolls back. An \textbf{unstable}
equilibrium is the top of a hill: the slightest nudge and it rolls away, never to
return. A \textbf{phase line} is the one-dimensional version of this picture, and
the sign of $f(y)$ tells you which way the marble rolls.
\end{intuitionbox}

\subsection{The Phase Line}

Because the equation is autonomous, we can draw everything on a single vertical
$y$-axis, the \emph{phase line}. Mark the equilibria (roots of $f$). Between them,
$f$ has a constant sign. Where $f(y)>0$, the solution increases (draw an upward
arrow); where $f(y)<0$, it decreases (downward arrow). An equilibrium with arrows
pointing \emph{toward} it from both sides is stable; arrows pointing \emph{away} is
unstable.

\begin{focusbox}
\textbf{The derivative test for stability.} Near an equilibrium $y^*$, linearize:
$f(y)\approx f'(y^*)(y-y^*)$ since $f(y^*)=0$. So a small displacement
$u=y-y^*$ obeys $\dv{u}{t}\approx f'(y^*)\,u$, plain exponential behavior. Therefore
\[
  f'(y^*) < 0 \ \Rightarrow\ \text{stable (perturbation decays)},\qquad
  f'(y^*) > 0 \ \Rightarrow\ \text{unstable (perturbation grows)}.
\]
The sign of the derivative of $f$ at the equilibrium decides its fate. This is the
one-dimensional ancestor of an enormous theory.
\end{focusbox}

\section{Logistic Growth}

Exponential growth $y'=ky$ holds only while resources are unlimited. No real
population grows forever; food, space, and predators eventually bite. The simplest
fix, proposed by the Belgian mathematician Pierre Fran\c{c}ois Verhulst in 1838, is
to let the per-capita growth rate \emph{decline} as the population approaches a
ceiling $M$, the \textbf{carrying capacity}:
\[
  \dv{P}{t} = r\,P\left(1 - \frac{P}{M}\right).
\]

\begin{intuitionbox}
Read the two factors. When $P$ is tiny, $1-P/M\approx 1$, and growth is nearly
exponential, $P'\approx rP$ --- plenty of room, breed freely. As $P$ climbs toward
$M$, the factor $1-P/M$ shrinks toward $0$, throttling the growth. At $P=M$ growth
stops entirely. The population is self-limiting: its own size applies the brakes.
\end{intuitionbox}

\subsection{The Phase Line of the Logistic}

This is autonomous with $f(P)=rP(1-P/M)$, a downward parabola in $P$ crossing zero
at $P=0$ and $P=M$. Between them $f>0$ (growth); below $0$ and above $M$, $f<0$. So:
\begin{itemize}
\item $P^*=0$ is \textbf{unstable}: $f'(0)=r>0$. A tiny population grows away from
extinction.
\item $P^*=M$ is \textbf{stable}: $f'(M)=-r<0$. Populations above or below $M$ are
drawn back to it.
\end{itemize}
Every positive starting population, no matter how small or how large, ends up at the
carrying capacity $M$. Figure~\ref{ch37:logistic-phaseline} draws this verdict as a
phase line, and Figure~\ref{ch37:logistic} shows the family of solution curves all
funneling toward $P=M$.

\begin{figure}[H]
\centering
\begin{tikzpicture}[>=stealth]
% vertical phase line for f(P) = r P (1 - P/M), roots at 0 and M
\draw[->, line width=1pt] (0,-1.2) -- (0,5.2) node[above] {$P$};
% equilibria
\fill[white, draw=trapred, line width=1.1pt] (0,0) circle (3pt);   % P*=0 unstable (open)
\fill[targetgreen] (0,4) circle (3pt);                             % P*=M stable (filled)
\node[right=8pt] at (0,0) {$P^*=0$ \ (unstable)};
\node[right=8pt] at (0,4) {$P^*=M$ \ (stable)};
% sign of f below 0: f<0, arrow points down (away from 0)
\draw[trapred, ->, line width=1.4pt] (0,-0.4) -- (0,-1.0);
% between 0 and M: f>0, arrows point up (toward M, away from 0)
\draw[targetgreen, ->, line width=1.4pt] (0,0.6) -- (0,1.6);
\draw[targetgreen, ->, line width=1.4pt] (0,2.2) -- (0,3.2);
% above M: f<0, arrow points down (toward M)
\draw[targetgreen, ->, line width=1.4pt] (0,4.9) -- (0,4.4);
\node[left=4pt, gray] at (0,-0.7) {$f<0$};
\node[left=4pt, gray] at (0,2.0) {$f>0$};
\node[left=4pt, gray] at (0,4.65) {$f<0$};
\end{tikzpicture}
\caption{The phase line of the logistic equation $\dv{P}{t}=rP(1-P/M)$. Arrows show
the direction of motion: away from the open circle at $P=0$ (unstable) and toward the
filled circle at $P=M$ (stable). Reading a one-dimensional picture like this tells us
the long-term fate of every solution without solving anything.}
\label{ch37:logistic-phaseline}
\end{figure}

\begin{figure}[H]
\centering
\begin{tikzpicture}
\begin{axis}[
  width=12cm, height=6.5cm,
  axis lines=middle,
  xlabel=$t$, ylabel=$P$,
  xmin=0, xmax=10, ymin=0, ymax=12,
  xtick={0,2,4,6,8,10}, ytick={0,5,10},
  domain=0:10, samples=120,
]
% carrying capacity M = 10, r = 0.9, several initial conditions
\addplot[gray, dashed, line width=0.9pt] {10};
\node[gray] at (axis cs:8.6,10.8) {$M=10$};
% P(t) = M / (1 + A e^{-r t}), A = (M-P0)/P0
\addplot[curveblue, line width=1.2pt] {10/(1 + (9.5)*exp(-0.9*x))};   % P0=1
\addplot[curveblue, line width=1.2pt] {10/(1 + (3)*exp(-0.9*x))};     % P0=2.5
\addplot[curveblue, line width=1.2pt] {10/(1 + (0.25)*exp(-0.9*x))};  % P0=8
\addplot[tangentred, line width=1.2pt] {10/(1 + (-0.18)*exp(-0.9*x))};% P0 above M
\end{axis}
\end{tikzpicture}
\caption{Logistic solutions for several initial populations. Small starts show the
S-shaped (sigmoid) rise: nearly exponential, then bending over at the inflection
near $P=M/2$, then leveling at the carrying capacity $M=10$. A start \emph{above}
$M$ (red) decays down to $M$. All roads lead to $M$.}
\label{ch37:logistic}
\end{figure}

\subsection{Solving the Logistic Equation}

It is separable, though the integral needs partial fractions. Separate:
\[
  \frac{\dd P}{P\,(1 - P/M)} = r\,\dd t.
\]
Split the left fraction. Write $\dfrac{1}{P(1-P/M)} = \dfrac{1}{P} +
\dfrac{1/M}{1-P/M}$ (you can verify by combining over a common denominator). Then
\[
  \int\!\left(\frac{1}{P} + \frac{1/M}{1-P/M}\right)\dd P
  = \ln P - \ln\!\left(1 - \frac{P}{M}\right) = \ln\frac{P}{\,1-P/M\,}.
\]
Integrating the right side gives $rt+C$, so
\[
  \ln\frac{P}{1-P/M} = rt + C
  \quad\Longrightarrow\quad
  \frac{P}{1-P/M} = A\,e^{rt}, \quad A=e^{C}.
\]
Solving algebraically for $P$ (multiply out and collect) yields the clean
\textbf{logistic function}
\[
  \boxed{\,P(t) = \frac{M}{1 + A\,e^{-rt}}\,},\qquad
  A = \frac{M - P_0}{P_0}.
\]
As $t\to\infty$, $e^{-rt}\to 0$ and $P\to M$, confirming the phase-line verdict. The
curve has a single inflection point where it switches from concave-up to
concave-down, located exactly at $P=M/2$, the moment of fastest growth.

\begin{walkbox}
Why is the fastest growth at $P=M/2$? The growth \emph{rate} is
$f(P)=rP(1-P/M)$, a downward parabola in $P$. Its maximum is at the vertex, halfway
between its roots $0$ and $M$, namely $P=M/2$. So the population adds members
fastest when it is exactly half-full. Before that, too few breeders; after that, too
little room. The sweet spot is the middle.
\end{walkbox}

\begin{applicationbox}
\textbf{A yeast culture, worked.} A flask of yeast follows the logistic law with
carrying capacity $M=1000$ cells and intrinsic rate $r=0.7\ \text{hr}^{-1}$,
starting from $P_0=50$ cells. Then $A=(1000-50)/50 = 19$, so
\[
  P(t) = \frac{1000}{1 + 19\,e^{-0.7t}}.
\]
How long until the culture is half full ($P=500$)? Half-full happens when the
denominator is $2$, i.e.\ $19e^{-0.7t}=1$:
\[
  t = \frac{\ln 19}{0.7} \approx \frac{2.944}{0.7}\approx 4.2\ \text{hours}.
\]
That is also the moment of fastest growth. Verhulst's 1845 fit of this curve to
Belgian census data was one of the first genuinely predictive uses of a
differential-equation model in the social sciences.
\end{applicationbox}

\begin{historybox}
Verhulst introduced the logistic curve in 1838 and named it \emph{logistique} in
1845, but his work was largely forgotten for eighty years. It was rediscovered
independently in the 1920s by Raymond Pearl and Lowell Reed, who fit it to the
United States population and (over-)confidently predicted an American ceiling of
about 197 million. Reality blew past that, which carries its own lesson: the
logistic model is a beautiful \emph{idealization}, exact for nothing, illuminating
for almost everything. A model's job is insight, not prophecy.
\end{historybox}

\section{Existence, Uniqueness, and the Limits of Formulas}

We have leaned on the comforting belief that an initial condition picks out exactly
one solution. When is that guaranteed?

\begin{formalbox}
\textbf{Existence and Uniqueness (Picard--Lindel\"of, stated plainly).} Consider the
initial value problem $\dv{y}{x}=f(x,y)$, $y(x_0)=y_0$. If $f$ and its partial
derivative $\pdv{f}{y}$ are continuous on some rectangle containing $(x_0,y_0)$, then
on some interval around $x_0$ there exists \emph{one and only one} solution passing
through $(x_0,y_0)$.
\end{formalbox}

This is the theorem that makes the slope-field picture trustworthy: continuous,
non-crossing solution curves filling the plane, exactly one through each point. When
the hypotheses fail, weird things can happen --- solutions can branch or blow up in
finite time.

\begin{trapbox}
The equation $\dv{y}{x}=\sqrt{\,\abs{y}\,}$ with $y(0)=0$ has \emph{two} solutions
through the origin: the constant $y\equiv 0$, and the curve $y=\tfrac14 x^2$ (for
$x\ge 0$). Uniqueness fails because $\pdv{f}{y} = \tfrac{1}{2\sqrt{y}}$ blows up at
$y=0$ --- the smoothness hypothesis is violated exactly there. The river forks. This
is not a pathology to fear but a reminder: uniqueness is a \emph{theorem with
hypotheses}, not a law of nature.
\end{trapbox}

\begin{aestheticbox}
Most first-order ODEs that arise in practice cannot be solved with a clean formula
at all --- $\dv{y}{x}=x^2+y^2$ has no elementary solution. Yet the slope field still
exists, the phase line still tells us about equilibria, and numerical methods (Euler
and its descendants) can trace any particular solution to whatever accuracy we
like. This is the quiet liberation of the subject: \emph{understanding} a
differential equation --- its long-term behavior, its equilibria, its qualitative
shape --- does not require a formula. The geometry comes first and goes furthest.
The formula, when it exists, is a bonus.
\end{aestheticbox}

\section{Exercises}

\begin{exercise}
Solve the separable equation $\dv{y}{x} = 3x^2 y$ with $y(0)=5$. Identify the lost
constant solution and check whether your general formula recovers it.
\end{exercise}

\begin{exercise}
Solve $\dv{y}{x} = \dfrac{1+y^2}{1+x^2}$. (Hint: both sides integrate to arctangents.)
Express the answer as an implicit relation, then solve for $y$ when $y(0)=1$.
\end{exercise}

\begin{exercise}
A sample of a radioactive isotope decays to $80\%$ of its initial mass in $3$ days.
Find its half-life, and determine how long until only $10\%$ remains.
\end{exercise}

\begin{exercise}
A turkey at $5^\circ$C is placed in a $175^\circ$C oven. After $30$ minutes its
internal temperature is $40^\circ$C. Using Newton's law of heating, find the
constant $k$ and predict when the turkey reaches $75^\circ$C.
\end{exercise}

\begin{exercise}
Solve the linear equation $\dv{y}{x} + \dfrac{1}{x}\,y = \cos x$ for $x>0$ by the
integrating-factor method. (Hint: $\mu = x$, and $\int x\cos x\,\dd x$ needs parts.)
\end{exercise}

\begin{exercise}
Solve $x\,\dv{y}{x} - 2y = x^3 e^{x}$. First put it in standard form, then find the
integrating factor. Show that $\mu = x^{-2}$ and finish the problem.
\end{exercise}

\begin{exercise}
For the autonomous equation $\dv{y}{t} = y^3 - y$, find all equilibria, draw the
phase line, and classify each equilibrium as stable or unstable using the sign of
$f'$ at each.
\end{exercise}

\begin{exercise}
For the autonomous equation $\dv{y}{t} = (y-1)(y-3)^2$, find the equilibria and
classify them. One of them is \emph{semi-stable} (attracting on one side, repelling
on the other) because $f$ touches zero without changing sign there. Identify which,
and explain using the phase-line arrows.
\end{exercise}

\begin{exercise}
A lake with carrying capacity $M=20{,}000$ fish currently holds $4{,}000$ fish and
grows logistically with $r=0.4\ \text{yr}^{-1}$. Write $P(t)$ explicitly, find when
the population reaches $10{,}000$, and find when it is growing fastest.
\end{exercise}

\begin{exercise}
\textbf{Mixing problem.} A tank holds $100$ liters of pure water. Brine containing
$0.5$ kg of salt per liter flows in at $4$ L/min, and the well-stirred mixture
flows out at the same rate. Let $S(t)$ be the kilograms of salt in the tank. Argue
that $\dv{S}{t} = 2 - \dfrac{4}{100}S$, solve this first-order linear equation with
$S(0)=0$, and find the long-run amount of salt. (Hint: the equilibrium is where
$S'=0$.)
\end{exercise}

\begin{exercise}
Show by direct substitution that the logistic function
$P(t)=\dfrac{M}{1+Ae^{-rt}}$ satisfies $\dv{P}{t}=rP(1-P/M)$. (This verifies the
solution without re-deriving it, a good habit: solving is hard, checking is easy.)
\end{exercise}

\begin{exercise}
The equation $\dv{y}{x} = x^2 + y^2$ has no elementary closed-form solution. Without
solving it, sketch its slope field for $-2\le x,y\le 2$ by computing the slope at a
handful of points, and describe qualitatively how a solution starting at $(0,0)$
behaves as $x$ increases. (This is the spirit of the final box: geometry without
formulas.)
\end{exercise}
